\chapter{Modelo e implementación}\label{cap3:modelo}

\section{Formulación del problema}

Se considera el dominio cuadrado adimensional $\Omega = [0,1]^2$, donde las
coordenadas han sido normalizadas por la longitud de onda del láser
($\lambda = 638\,\mathrm{nm}$, láser de diodo rojo). Bajo esta normalización,
el número de onda adimensional es $\tilde{k} = 2\pi$.

Se plantean dos experimentos de complejidad creciente (completados), con dos adicionales planificados como trabajo futuro:

\subsection{Experimento 1 (NB01): Validación 1D}

La ecuación de Helmholtz unidimensional en $x \in [0,1]$ con condiciones de
Dirichlet:
%
\begin{equation}
    \frac{d^2 E}{dx^2} + k^2 E = 0, \quad
    E(0) = 1, \quad E(1) = \cos(k),
    \label{eq:helmholtz1d}
\end{equation}
%
admite la solución analítica $E_\text{exacta}(x) = \cos(kx)$, utilizada para
validar la arquitectura y la función de pérdida.

\subsection{Experimento 2 (NB02): Validación 2D con campo complejo}

La ecuación de Helmholtz 2D en $\Omega = [0,1]^2$ con condición de onda plana
diagonal en los cuatro bordes:
%
\begin{equation}
    E_\text{exacta}(x,y) = e^{i(k_x x + k_y y)}, \quad
    k_x = k_y = \frac{k}{\sqrt{2}},
    \label{eq:solucion2d}
\end{equation}
%
que satisface $k_x^2 + k_y^2 = k^2$. La condición de frontera se impone
separando la solución en sus partes real e imaginaria:
%
\begin{equation}
    E_\text{real}(x,y)\big|_{\partial\Omega} = \cos(k_x x + k_y y), \qquad
    E_\text{imag}(x,y)\big|_{\partial\Omega} = \sin(k_x x + k_y y).
\end{equation}

% ──────────────────────────────────────────────────────────────────────────────
\section{Arquitectura SIREN}

Se emplea una arquitectura SIREN con los parámetros consolidados a partir de la
validación 1D (NB01) y escalados para el caso 2D (NB02):

\begin{itemize}
    \item \textbf{Entrada:} coordenadas espaciales $(x, y) \in [0,1]^2$ (o
    $(x) \in [0,1]$ para 1D).
    \item \textbf{Capas ocultas:} $L = 5$ capas de dimensión $d$ neuronas cada
    una ($d = 64$ en NB01, $d = 128$ en NB02).
    \item \textbf{Activación:} $\phi(z) = \sin(\omega_0 z)$, con $\omega_0 = 1.0$.
    \item \textbf{Salida:} capa lineal de dos neuronas $(E_\text{real}, E_\text{imag})$
    para el caso 2D, o una neurona $E$ para el caso 1D.
    \item \textbf{Parámetros totales:} ${\sim}8{,}400$ (NB01) y $66{,}690$ (NB02).
\end{itemize}

La elección $\omega_0 = 1.0$, en lugar del valor $\omega_0 = 30$ recomendado
por \textcite{sitzmann2020siren} para señales de imagen, está justificada por la
naturaleza del problema: el dominio adimensional $[0,1]^2$ con $k = 2\pi$ no
requiere capturar frecuencias ultra-altas, y valores mayores de $\omega_0$
desestabilizan el entrenamiento con L-BFGS.

La inicialización sigue el esquema de \textcite{sitzmann2020siren}:
$\mathbf{W}_0 \sim \mathcal{U}(-1/n_\text{in}, 1/n_\text{in})$ para la primera
capa y $\mathbf{W}_\ell \sim \mathcal{U}(-\sqrt{6/d}, \sqrt{6/d})$ para capas
intermedias, garantizando que las distribuciones de activación sean
aproximadamente uniformes en $[-1, 1]$.

% ──────────────────────────────────────────────────────────────────────────────
\section{Función de pérdida y calibración de pesos}

La función de pérdida total para el caso 2D es:
%
\begin{equation}
    \mathcal{L} = \mathcal{L}_\text{datos}
    + \lambda_\text{fís}\!\left(
    \mathcal{L}_{R_\text{real}} + \mathcal{L}_{R_\text{imag}}
    \right),
    \label{eq:loss_2d}
\end{equation}
%
con peso de física $\lambda_\text{fís} = 0.1$. Este valor fue calibrado
empíricamente mediante un experimento de ablación (véase la
Sección~\ref{sec:ablacion}): valores superiores ($\lambda = 1.0$) sobrepesan la
física e impiden la convergencia en 2D, dado que los dos residuos (real e
imaginario) duplican efectivamente la contribución física; valores inferiores
($\lambda = 0.01$) producen campos que violan moderadamente la ecuación de
Helmholtz.

Los residuos de Helmholtz se calculan mediante diferenciación automática de
segundo orden \parencite{baydin2018autodiff}:
%
\begin{equation}
    R_\text{real}(\mathbf{r}) = \nabla^2 E_\text{real}(\mathbf{r})
    + k^2 E_\text{real}(\mathbf{r}), \qquad
    R_\text{imag}(\mathbf{r}) = \nabla^2 E_\text{imag}(\mathbf{r})
    + k^2 E_\text{imag}(\mathbf{r}),
\end{equation}
%
donde el Laplaciano se obtiene con \texttt{torch.autograd.grad} activando la
opción \texttt{create\_graph=True} para permitir diferenciación de segundo orden.
Es crítico que $k^2$ se convierta a tipo \texttt{float} de Python antes de
multiplicar por el tensor de la red, para evitar conflictos de tipo entre
\texttt{numpy.float64} y \texttt{torch.float32}.

% ──────────────────────────────────────────────────────────────────────────────
\section{Estrategia de muestreo}

\subsection{Puntos interiores (colocación)}

Para el caso 2D, los $N_c = 3{,}000$ puntos de colocación interiores se generan
mediante LHS \parencite{mckay1979lhs} usando \texttt{scipy.stats.qmc.LatinHypercube}
con dimensión $d = 2$. LHS divide cada dimensión en $N_c$ estratos iguales de
longitud $1/N_c$ y coloca exactamente un punto en cada estrato, garantizando
cobertura uniforme del dominio sin los patrones redundantes de una malla
cartesiana.

\subsection{Puntos de frontera}

Los $N_b = 300$ puntos de frontera por borde se distribuyen uniformemente sobre
cada uno de los cuatro lados del dominio $\partial\Omega$. Para cada lado, se
evalúa la solución exacta y se
construye el par $(\mathbf{r}^b_j,\, E(\mathbf{r}^b_j))$ utilizado en
$\mathcal{L}_\text{datos}$.

% ──────────────────────────────────────────────────────────────────────────────
\section{Estrategia de optimización}

El entrenamiento sigue la estrategia bifásica descrita en el
Capítulo~\ref{cap2:marcos}, con los hiperparámetros consolidados en la
Tabla~\ref{tab:hiperparametros}.

\begin{table}[htbp]
\centering
\begin{threeparttable}
\caption{Hiperparámetros de entrenamiento consolidados.}
\label{tab:hiperparametros}
\begin{tabular}{lll}
\toprule
Hiperparámetro & Valor & Aplicación \\
\midrule
Tasa de aprendizaje Adam ($\eta$) & $10^{-3}$ & NB01, NB02 \\
Épocas máximas Adam             & 15{,}000  & NB01, NB02 \\
Paciencia (\textit{early stopping}) & 800 épocas & NB02 \\
$\delta$ mejora mínima          & $10^{-7}$ & NB02 \\
Iteraciones máximas L-BFGS      & 500 (NB01), 1{,}000 (NB02) & — \\
Historial L-BFGS                & 100       & NB01, NB02 \\
Semilla global (SEED)           & 42        & NB01, NB02 \\
Peso de física $\lambda_\text{fís}$ & 0.1  & NB02 \\
\bottomrule
\end{tabular}
\begin{tablenotes}\small
    \item Todos los experimentos se ejecutaron en GPU NVIDIA RTX 5050
    (CUDA 12.6, PyTorch 2.4). La semilla garantiza reproducibilidad.
\end{tablenotes}
\end{threeparttable}
\end{table}

% ──────────────────────────────────────────────────────────────────────────────
\section{Métricas de validación}

\subsection{Error $L^2$ relativo (NB01 y NB02)}

El error $L^2$ relativo mide la discrepancia global entre la predicción de la
red $\hat{E}$ y la solución exacta $E_\text{exacta}$, evaluado sobre una malla
uniforme de $100 \times 100$ puntos:
%
\begin{equation}
    \varepsilon_{L^2} = \frac{\|\hat{E} - E_\text{exacta}\|_2}
    {\|E_\text{exacta}\|_2} \times 100\%.
    \label{eq:l2}
\end{equation}

El umbral de la tesis es $\varepsilon_{L^2} < 5\%$.

% ──────────────────────────────────────────────────────────────────────────────
\section{Implementación modular}

El código fuente está organizado en módulos reutilizables en \texttt{src/}:

\begin{description}
    \item[\texttt{models.py}] Clase \texttt{PINN\_2D\_SIREN} que implementa
    la arquitectura SIREN con la inicialización de \textcite{sitzmann2020siren}.
    El método \texttt{forward(x)} aplica la activación sinusoidal capa por capa
    y retorna el campo complejo como un tensor de dos columnas
    $(E_\text{real}, E_\text{imag})$.

    \item[\texttt{losses.py}] Función \texttt{pinn\_loss\_2d(model, xy\_int,
    xy\_bc, E\_bc, k, lam)} que calcula la pérdida total~\eqref{eq:loss_2d}.
    El cálculo del Laplaciano utiliza \texttt{torch.autograd.grad} con
    \texttt{create\_graph=True} para habilitar la diferenciación de segundo orden.
    El escalar $k$ se convierte explícitamente a \texttt{float(k)} antes de
    multiplicar por el tensor de la red, evitando conflictos de dtype entre
    \texttt{numpy.float64} y \texttt{torch.float32}.

    \item[\texttt{training.py}] Loop de entrenamiento que integra las fases
    Adam y L-BFGS con \textit{early stopping}.

    \item[\texttt{utils.py}] Métricas ($L^2$, $R^2$, RMSE, MAE), funciones de
    muestreo LHS y rutinas de visualización.
\end{description}

Los experimentos están reproducidos en los notebooks
\texttt{01\_pinn\_helmholtz\_1d\_validation.ipynb} y
\texttt{02\_pinn\_helmholtz\_2d\_complex\_field.ipynb}.
Los resultados numéricos están archivados en formato JSON en \texttt{results/}.
